\chapter{Little and Great Tradition}

\section{Tradition: Meaning and Forms}

Before studying the modernisation of Indian tradition, one must understand what \textbf{tradition} itself is.

\begin{KeyConcept}
\begin{itemize}
\item \textbf{Tradition} --- any custom, practice or ritual that is forwarded from one generation to the next, either \textbf{orally} or in \textbf{written form}.
\item If it continues for generations together, it becomes established as a tradition.
\item The key requirement of a tradition is \textbf{longevity} --- it must have been sustained over a long time: through upheavals, earthquakes and other disturbances, the practice still survives. That is what makes it a tradition.
\end{itemize}
\end{KeyConcept}

\subsection{Oral vs written tradition}

\begin{itemize}
\item \textbf{Written tradition} is generally regarded as mainstream, elitist, superior, authentic and original.
\item \textbf{Oral tradition} is regarded as not mainstream, not classical, not elitist --- local, limited to a few people and not popular.
\begin{itemize}
\item Written things become popular easily because they are printed and copies are distributed cheaply; written things have reach, which oral ones lack (e.g. a statement written on social media spreads, while a spoken remark reaches only one friend).
\end{itemize}
\end{itemize}

\section{Robert Redfield and the Study of Tradition}

The systematic study of tradition was initiated by an American anthropologist, \textbf{Robert Redfield}.

\begin{DataFact}
\begin{itemize}
\item Robert Redfield --- American anthropologist.
\item Conducted a field study in \textbf{Tepoztlán}, a village in Mexico.
\item There he identified \textbf{two types of tradition}: \textbf{Little Tradition} and \textbf{Great Tradition}.
\end{itemize}
\end{DataFact}

In every society --- village or city --- every civilisation has its own traditions and customs, practised locally (orally or in writing) but absent in other pockets. Some traditions are practised locally, others universally.

\begin{KeyConcept}
\begin{itemize}
\item \textbf{Great Tradition} is always the \textbf{written} one.
\item \textbf{Little Tradition} is always the \textbf{oral} one.
\item The little tradition is practised at the local/village level; the great tradition is known universally, all over India.
\end{itemize}
\end{KeyConcept}

\section{Kishan Garhi Village and Milton Singer's Study}

Redfield's students visited Indian villages to observe these traditions, just as Redfield had studied Tepoztlán. One such village was \textbf{Kishan Garhi} in Uttar Pradesh, studied by \textbf{Milton Singer} (a student of Redfield).

\begin{ExampleBox}
\begin{itemize}
\item \textbf{Govardhan} --- the local practice at Kishan Garhi. It is believed that some 40 km from the village, Lord Krishna descended to help villagers protect their cattle from a heavy rain by lifting the entire \textbf{Govardhan} mountain on his little finger (while playing the flute on the other hand).
\item Today a local practice called \textbf{Govardhan} is celebrated around this: cattle dung is treated as wealth (a source of methane), and rituals are conducted around Govardhan because it sustains the local economy.
\item This is a local practice --- the whole of India cannot be expected to celebrate Govardhan or have a Govardhan holiday --- so it remains a \textbf{little tradition}.
\item The same story is used by the Yadav people to trace their descent/lineage to Lord Krishna (Yadu Maharaj). The story is passed \textbf{orally}, so it is a little tradition.
\end{itemize}
\end{ExampleBox}

The terms 'little tradition' and 'great tradition' were coined by Robert Redfield; the examples studied in India come from Redfield's students.

\section{Universalization of Little Tradition}

It is not true that little traditions always remain little. Over a period of time they may become great, and once great, the oral little tradition may even acquire a written form to justify its greatness.

\begin{ExampleBox}
\begin{itemize}
\item \textbf{Sannauli $\rightarrow$ Raksha Bandhan}. In Kishan Garhi there is a local custom called \textbf{Sannauli}: a brother must visit his married sister once a year, carrying gifts for the priest who arranged the marriage and for the sister's in-laws; the sister applies a tilak (with barley grains) on his forehead and ties a thread (dhaga) on his wrist as a mark of the visit.
\item This local practice is celebrated today in the name of \textbf{Raksha Bandhan} --- it has become a great tradition. Raksha Bandhan was never originally a great tradition; its 'little' form was Sannauli.
\item Today Raksha Bandhan has a national holiday and written texts to justify it (its history is found even on Wikipedia); nothing is written about Sannauli, which remains oral.
\end{itemize}
\end{ExampleBox}

\begin{itemize}
\item When a little tradition becomes great or universalised --- celebrated all over India in different forms --- the process is called \textbf{universalization of little tradition}.
\item The terms \textbf{universalization} and \textbf{parochialization} were coined by \textbf{McKim Marriott}, another student of Redfield.
\end{itemize}

\begin{ExampleBox}
\begin{itemize}
\item \textbf{Karwa Chauth} --- today celebrated all over India, but it began as a local custom of western Rajasthan and Gujarat and then took universalised form. It was never originally an all-India (great) tradition; it was oral, and even today nothing is written about it, yet it is universally practised and has become part of the great tradition.
\item \textbf{Jaya $\rightarrow$ Bharata $\rightarrow$ Mahabharata}. The great epic Mahabharata is a great tradition, but its original form was \textbf{Jaya} --- a local text of about 3,000 verses --- which grew through Bharata into the Mahabharata (the Sanskrit epic attributed to Vyasa, who in reality wrote only a few parts).
\end{itemize}
\end{ExampleBox}

\section{Parochialization of Great Tradition}

\begin{itemize}
\item When a great tradition becomes little --- takes a localised form in different villages --- the process is called \textbf{parochialization}.
\end{itemize}

\begin{ExampleBox}
\begin{itemize}
\item \textbf{Govardhan} is a parochialised form of the Govardhan event (Krishna saving the cattle). The full Krishna narrative --- his childhood story and life --- is known all over India (great tradition); in different villages it has taken a parochialised, local form.
\item \textbf{Diwali} was initially a local festival celebrating the welcome of Rama (on his return after killing Ravana), but has become part of the great tradition --- popularised through the \textbf{Ramayana} and \textbf{Ramlila}. Ramlila was itself practised locally, and every village began projecting it to spread awareness of the great heroes among the masses.
\item \textbf{Harikatha} (the stories of Krishna) was practised locally but gradually took a great avatar and became part of the great tradition.
\item \textbf{G. S. Ghurye} also discussed this: local customs limited to the Hindu Vedic belt diffused all over India with the help of priests and customs like \textbf{Ashwamedha} --- little traditions that became great once practised everywhere.
\end{itemize}
\end{ExampleBox}

\begin{KeyConcept}
\begin{itemize}
\item \textbf{McKim Marriott} believes \textbf{social change in India can be explained with the help of the interaction between little tradition and great tradition}.
\item Just as \textbf{westernization} and \textbf{Sanskritization} are processes of social change, \textbf{universalization} (a local practice becoming universal) and \textbf{parochialization} (a universal practice taking a local avatar) are also processes of social change.
\end{itemize}
\end{KeyConcept}

\section{Examples from Kishan Garhi: Saurati and Naurata}

Milton Singer studied Kishan Garhi and its little traditions; McKim Marriott re-studied the same village and developed the concepts of universalization and parochialization. Kishan Garhi was treated as representative of all Indian villages because it was imbued with folk customs and local traditions (just as Gandhi saw the village as representing India).

\begin{ExampleBox}
\begin{itemize}
\item \textbf{Festival of Lights} --- celebrated around a local goddess \textbf{Saurati} (also written Saurata; S-A-U-R-T-I / S-A-U-R-T-H-A), a symbol of wealth and prosperity.
\item Marriott believed \textbf{Saurati is a parochialised form of Lakshmi} (Lakshmi is the all-India Sanskritic Hindu goddess of wealth and prosperity, who at the local level becomes Saurati).
\item \textbf{Naurata} --- a local goddess in Kishan Garhi worshipped for \textbf{nine days/nights} (made of mud). Marriott believed \textbf{Naurata became Navratri all over India} --- the origin of Navratri lies in Naurata, and the nine-night Navratri customs (fasting, food restrictions) practised all over India originated in that village.
\item \textbf{Naurata has also influenced Durga Puja and Kali Puja} (per Marriott).
\end{itemize}
\end{ExampleBox}

\begin{CriticalPoint}
\begin{itemize}
\item The teacher flagged that the exact direction of the deity examples is easy to confuse, and instructed students to follow the written notes. The written version is:
\item \textbf{Parochialization} --- a process of \textbf{localization} of the great tradition: in Kishan Garhi, the mud female deity \textbf{Naurata} is worshipped for nine consecutive nights; Marriott believes \textbf{Durga has been parochialised into Naurata} (the name itself deriving from 'Navratri').
\item \textbf{Universalization} --- in the Festival of Lights, the local goddess \textbf{Saurati} would have been \textbf{universalised into the goddess Lakshmi} of the great tradition.
\end{itemize}
\end{CriticalPoint}

In conclusion, seen through festivals and deities, the religion of the village of Kishan Garhi may be said to have originated from a continuing process of communication between the little and the great tradition.

\section{Sanskritization as Parochialization}

\begin{ComparisonBox}
\begin{itemize}
\item \textbf{All-India Sanskritic Hinduism} = the \textbf{great tradition} --- found everywhere across India; there is no state without a temple of the Hindu gods.
\item \textbf{Local Hinduism} (local gods like \textbf{Karni Mata, Mansa Devi, Murugan}) = the \textbf{little tradition} --- practised locally, limited to a specific region, by few people, not found elsewhere.
\item Great tradition is largely among literate, Sanskrit-knowing elites; little tradition is oral, among the illiterate, at mass level but local.
\end{itemize}
\end{ComparisonBox}

\begin{KeyConcept}
\begin{itemize}
\item \textbf{Sanskritization} is the process by which local Hinduism borrows from all-India Sanskritic Hinduism --- i.e. low-caste, tribal or other groups emulate the \textbf{twice-born (dvija)} upper castes, their lifestyle, customs and even their gods.
\item Therefore \textbf{Sanskritization is nothing but the parochialization of all-India Sanskritic Hinduism}.
\item Sanskritization can also be called \textbf{traditionalization} --- the emulation of dvija traditions by the non-dvija low castes.
\end{itemize}
\end{KeyConcept}

\begin{ExampleBox}
\begin{itemize}
\item \textbf{Coorg (Kodagu)} example (from M. N. Srinivas): the Coorgs were invaded by the \textbf{Lingayats} (Shaivites, worshippers of Shiva). The Coorgs then claim \textbf{Kshatriya status}; Shiva (a great-tradition deity) got parochialised among the Coorg.
\item The Coorgs began to consider the \textbf{Kaveri} as a sacred river (Kaveri is a sacred goddess among Hindus), and even wrote a \textbf{Kaveri Purana} --- a local Purana that does not exist among the great Puranas.
\item This shows great tradition getting localised/parochialised among a local group.
\end{itemize}
\end{ExampleBox}

\section{Great \& Little Tradition and Social Change: Yogendra Singh's Framework}

The great and little traditions form part of our culture, and they must be understood to grasp \textbf{Yogendra Singh's} book \textbf{Modernization of Indian Tradition}. There is no single 'Indian tradition' --- traditions vary from village to village and from the local to the great level.

\begin{KeyConcept}
\begin{itemize}
\item Tradition does not die; it \textbf{adapts}. When exposed to modern/colonial forces, a tradition like Holi adapts rather than disappears --- today Holi is played with modern water jets and AK-47-shaped pichkaris; Diwali uses environment-friendly, carbon-free, sustainable firecrackers.
\item This is the unique feature of tradition --- it is continuous and passed from generation to generation.
\end{itemize}
\end{KeyConcept}

Yogendra Singh studies \textbf{modernization} as the change introduced by colonization (modern education, modern laws, parliamentary structure, democracy, fundamental rights --- all came from outside) and its interaction with existing traditions. He looks at changes in \textbf{every realm}, at both the \textbf{cultural level} and the \textbf{structural level}.

\begin{ComparisonBox}
\begin{itemize}
\item Yogendra Singh's framework (the table to remember):
\item \textbf{Sources of change} --- cultural structure vs social structure.
\item \textbf{Cultural structure} --- \textbf{little tradition} and \textbf{great tradition}.
\item \textbf{Social structure} --- \textbf{micro-structure} and \textbf{macro-structure}.
\item \textbf{Direction of change} --- \textbf{orthogenetic} (change from within) and \textbf{heterogenetic} (change from outside).
\item Examples: Naurata $\rightarrow$ Durga (Durga is a universalised form of Naurata) is change \textbf{from within} = orthogenetic; \textbf{Sanskritization} is orthogenetic (from within India); \textbf{westernization} is heterogenetic (from outside).
\end{itemize}
\end{ComparisonBox}

\begin{itemize}
\item \textbf{Orthogenetic change} (from within): \textbf{Sanskritization / traditionalization}.
\item \textbf{Heterogenetic change} (from outside): \textbf{primary westernization}, \textbf{Islamization} (and Sufi culture --- Sufi songs/'Sufization' of Bollywood coming from outside). Secondary westernization is largely heterogenetic but can also be orthogenetic.
\begin{itemize}
\item Primary westernization is heterogenetic; secondary (and tertiary) westernization are orthogenetic in form.
\item The primary people (upper-caste Brahmins and Banias, especially Brahmins) got western education first, and were the first to be emulated by the lower castes --- who in emulating them also emulated western culture; hence secondary westernization faced some cultural resistance.
\end{itemize}
\end{itemize}

\begin{KeyConcept}
\begin{itemize}
\item \textbf{Islamization} began with the advent of the Sultanate: it introduced things alien to Indian civilisation --- new land systems, and a warrior class with its own army (the first such class in India).
\item With the Sultanate and then the Mughals, a harmony and peaceful mutual coexistence was gradually established between the local indigenous (Vedic) people and the Muslim outsiders --- though there were also periods of war (which Yogendra Singh does not dwell on; he avoids the conflict dimension and is not influenced by Marxism here).
\end{itemize}
\end{KeyConcept}

\begin{KeyConcept}
\begin{itemize}
\item \textbf{Yogendra Singh's definition of modernization}: an \textbf{ongoing process} from a traditional, agrarian Indian society to a modern, industrial (Western) society.
\item He insists modernization in India should \textbf{not be seen as a binary opposition} --- in the West modernization came at the cost of tradition, but in India traditions get modernised; tradition and modernity \textbf{coexist}.
\item Example: \textbf{caste} is a traditional system, yet it has taken a modern avatar through \textbf{caste associations} (e.g. an all-India Thakur association) --- modern democratic associations formed around traditional identities, possible only because of the Constitution.
\item Example: \textbf{Panchayati Raj Institution} is the democratic (constitutionalised) avatar of the old local \textbf{caste panchayat} --- the panchayat did not end, it got modernised/institutionalised.
\item Example: worship of \textbf{Lord Jagannath} now happens digitally --- pilgrimage, digital aarti, online donation --- the tradition persists in a modernised avatar.
\end{itemize}
\end{KeyConcept}

Yogendra Singh's theory is considered the \textbf{most comprehensive} because it covers every level: changes in little tradition and great tradition, and changes at micro-structure and macro-structure.

\begin{itemize}
\item \textbf{Macro-structure}: modern bureaucracy, parliament, meritocracy --- systems India did not have before (recruitment was patronised/nepotistic, not a rational written process). This is a macro-level change (rational bureaucracy, open competitive exam with universalised treatment).
\item \textbf{Micro-structure}: individuals --- e.g. the traditional role of women (housewife) and of Dalits has changed; women now work at par with men (gender equality / women empowerment). This micro-level change also affects the macro-structure of the family (families weakened/separated, 'living apart together' arrangements).
\end{itemize}

\begin{DataFact}
\begin{itemize}
\item The \textbf{modernization theory} was originally given by classical theorists --- \textbf{Comte, Marx, Weber, Parsons} --- but it explained the modernising process \textbf{of the West}.
\item In contemporary times, especially the \textbf{post-Cold War} period, modernization theories explain the modernization \textbf{of the non-West} --- every country can be placed within \textbf{Rostow's five stages}: traditional; pre-take-off; take-off; drive to maturity; and the age of mass consumption.
\item Yogendra Singh explains modernization through the process of \textbf{colonization}, but he does not limit himself to colonization --- he also studies the impact of \textbf{westernization}, \textbf{Islamization} and \textbf{Sanskritization}.
\end{itemize}
\end{DataFact}

The detailed explanation of the framework is continued in the next lecture (Modernization of Indian Tradition).

\begin{QuickRevision}
\begin{itemize}
\item Tradition = custom/practice/ritual passed down generations (orally or written), sustained by longevity.
\item Written tradition = mainstream/elitist/superior; oral tradition = local/limited/not popular.
\item Robert Redfield (American anthropologist) studied Tepoztlán, Mexico and coined \textbf{Little Tradition} (oral, local) and \textbf{Great Tradition} (written, universal).
\item Milton Singer studied Kishan Garhi (UP); McKim Marriott re-studied it and coined \textbf{universalization} (little $\rightarrow$ great) and \textbf{parochialization} (great $\rightarrow$ little).
\item Examples: Sannauli $\rightarrow$ Raksha Bandhan; Karwa Chauth (local Rajasthan/Gujarat $\rightarrow$ all-India); Jaya $\rightarrow$ Bharata $\rightarrow$ Mahabharata; Diwali + Ramlila; Govardhan (parochialised).
\item Kishan Garhi deities: Saurati (Festival of Lights) $\leftrightarrow$ Lakshmi; Naurata (9 nights) $\rightarrow$ Navratri, also influencing Durga Puja/Kali Puja.
\item Sanskritization = parochialization of all-India Sanskritic Hinduism; Coorg/Lingayat/Kaveri Purana example.
\item Yogendra Singh's framework: cultural (little/great) + structural (micro/macro) + orthogenetic (within) vs heterogenetic (outside); modernization = agrarian $\rightarrow$ industrial, tradition \& modernity coexist (not binary).
\end{itemize}
\end{QuickRevision}

